\pagestyle{plain}

\chapter{CAPÍTULO 3: Meta-análise}


\begin{flushright}
Catuxe Varjão de Santana Oliveira  \footnote[1]{Programa de Pós-graduação em Ciências da Propriedade Intelectual-PPGPI/Universidade Federal de Sergipe-UFS. *E-mail:  \href{mailto:catuxe@academico.ufs.br}{\color{blue}{\underline{catuxe@academico.ufs.br}}}}
\\[0.3cm]
Luiz Diego Vidal Santos \footnote[2]{Universidade Estadual de Feira de Santana - UEFS, Brasil. E-mail: \href{mailto:ldvsantos@uefs.br}{\color{blue}{\underline{ldvsantos@uefs.br}}}}
\\[0.3cm]
Paulo Roberto Gagliardi \footnotemark[1]
\end{flushright}


% resumo em português — bookmark suprimido para manter hierarquia do capítulo
{\renewcommand{\PRIVATEbookmarkthis}[1]{}%
\begin{resumoumacoluna}
A assimetria entre a riqueza biocultural dos Sistemas Agroecológicos Tradicionais (SSAT) e sua invisibilidade econômica configura um impasse epistemológico e metodológico que impede a tradução desses saberes tácitos em ativos intangíveis auditáveis, perpetuando a vulnerabilidade territorial e econômica das comunidades quilombolas diante de cadeias produtivas globais. Esta investigação tem como objetivo identificar e analisar o potencial dos sistemas agroecológicos tradicionais como ativos intangíveis, mapeando dimensões de valor cultural, científico e econômico que orientem sua valoração comunitária e incorporação em estratégias de gestão do conhecimento. Aplicada a comunidades quilombolas do semiárido brasileiro, a pesquisa integrou instrumentos psicométricos para mensuração de capital cognitivo-cultural, modelos econômicos baseados no framework VAIC e validação comunitária participativa atendendo aos critérios VRIN (Valiosos, Raros, Inimitáveis e Não-substituíveis), fornecendo a base probatória necessária para segurança fundiária, negociação equitativa em cadeias de valor da bioeconomia e fortalecimento da soberania epistêmica comunitária. Os resultados demonstram que a integração psicométrica-econômica-participativa não apenas quantifica o valor estratégico dos conhecimentos tradicionais, mas também instrumentaliza as comunidades para protagonismo ativo na gestão de seus ativos bioculturais, estabelecendo um protocolo replicável para a transição de saberes tácitos a ativos econômicos tangíveis em contextos de justiça cognitiva e repartição equitativa de benefícios.

\vspace{\onelineskip}
\noindent
\textbf{Palavras-chave}: Valoração de ativos intangíveis; Psicometria aplicada; Bioeconomia; Soberania epistêmica; Sistemas agroecológicos tradicionais.
\end{resumoumacoluna}}

\newpage

\section{Introdução}

A economia do conhecimento contemporânea fundamenta-se na premissa de que ativos intangíveis constituem o núcleo determinante da competitividade organizacional e da sustentabilidade de sistemas produtivos, conforme demonstrado pelas teorias de valoração de capital intelectual propostas por \citeonline{Edvinsson1997} e operacionalizadas pelo modelo Value Added Intellectual Coefficient (VAIC) de \citeonline{Pulic2000}. Contudo, a aplicação dessas métricas permanece restrita ao ambiente corporativo formal, incapaz de capturar a complexidade multidimensional dos saberes tradicionais que, embora atendam rigorosamente aos critérios de recursos VRIN postulados por Barney \cite{Barney1991}, permanecem economicamente invisíveis devido à ausência de instrumentos e modelos de valoração adaptados à sua natureza tácita e distribuída. Essa lacuna metodológica configura o que denominamos hiato gesrencial-econômico, responsável pela subvalorização sistemática dos Saberes e Sistemas Agrícolas Tradicionais (SSAT) nos mercados de conhecimento e nas estratégias de gestão territorial das comunidades.

Segundo \citeonline{deCarvalho2016}, a invisibilidade econômica dos conhecimentos tradicionais não deriva de uma suposta inferioridade tecnológica ou de relevância marginal para a inovação contemporânea, mas resulta de uma assimetria fundamental entre os modos de codificação do saber. Tais saberes permanecem economicamente e politicamente invisibilizados não por falta de sofisticação, mas porque operam com outras ontologias, linguagens e regimes de registro, que colidem com os modos hegemônicos de codificação do conhecimento científico e econômico no Antropoceno \cite{deCarvalho2016,Agrawal2002}.

Enquanto a economia da inovação desenvolveu instrumentos sofisticados para mensurar ativos explícitos mediante indicadores de Pesquisa e Desenvolvimento (P\&D), patentes registradas e publicações indexadas, os conhecimentos tradicionais operam sob a lógica do conhecimento tácito descrita por \citeonline{Polanyi1966}, transmitidos oralmente e incorporados em práticas corporais e rotinas comunitárias que resistem à formalização cartesiana. Essa divergência epistemológica cria uma barreira de incomensurabilidade que impede a tradução direta dos saberes agroecológicos para as linguagens da contabilidade de ativos intangíveis \cite{Bastida2022}, resultando na exclusão dessas comunidades dos circuitos de valorização econômica da bioeconomia e dos mercados de serviços ecossistêmicos, apesar de deterem portfólios tecnológicos de adaptação climática e conservação da agrobiodiversidade comprovadamente superiores aos modelos industriais hegemônicos \cite{Laborde2023}.

A superação desse hiato transcende a simples tradução metodológica, exigindo uma arquitetura de codificação capaz de converter a complexidade fenomenológica dos saberes tradicionais em ativos epistêmicos auditáveis \cite{Boyd2020}. Nesse contexto, análises psicometrica aplicadas reconfigura-se como uma tecnologia de estruturação de dados, operando como a interface crítica entre a oralidade difusa e o rigor algorítmico \cite{McCrae2019}. Ao transmutar narrativas e práticas tácitas em vetores multidimensionais calibrados, a modelagem psicométrica gera a matéria-prima estruturada indispensável para o processamento por inteligência artificial \cite{Bengio2021}. Essa vetorização do conhecimento estabelece uma base probatória de alta fidelidade que captura a variância real das especializações em manejo e adaptação climática, superando os vieses que historicamente desqualificaram as tecnologias sociais e preparando o terreno para que modelos de \textit{Machine Learning} possam decodificar padrões complexos de valor \cite{Valdez2022}.

Entretanto, a vetorização de dados cognitivos não é um fim em si mesma, mas a infraestrutura necessária para converter a riqueza biocultural em ativos auditáveis. O Machine Learning atua aqui como decodificador de complexidade, correlacionando a proficiência agroecológica (o saber) com o Valor Adicionado (o retorno econômico, via VAIC). O objetivo não é apenas treinar um algoritmo, mas criar uma métrica de valor que seja legível pelo mercado (bioeconomia) e, ao mesmo tempo, fiel à realidade comunitária \cite{HosseiniyanKhatibip}. 

A vetorização dos saberes tradicionais transcende a mera engenharia de dados, ela constitui a infraestrutura crítica para a materialização econômica de ativos historicamente invisibilizados \cite{Hasani2020}. Nesta arquitetura, a Inteligência Artificial não opera apenas como ferramenta estatística, mas como um mecanismo de tradução ontológica capaz de correlacionar a complexidade sistêmica dos SSAT com indicadores de desempenho econômico. O framework VAIC de Pulic \cite{Pulic2000IC}, ao definir as métricas de valor agregado, fornece a gramática econômica necessária para que os algoritmos identifiquem padrões de rentabilidade no capital intangível comunitário, efetivando a transição de 'saberes locais' para 'ativos estratégicos' disputáveis em cadeias globais de valor.

Entretanto, a legitimidade dessa tradução econômica depende imperativamente de salvaguardas epistemológicas que impeçam a reprodução da violência colonial descrita por Santos \cite{Santos2007}. A validação comunitária, fundamentada na validade consequencial de Messick \cite{Messick1995} e na pesquisa-ação de Fals Borda \cite{FalsBorda1991}, não é apenas uma etapa de calibração do modelo, mas o dispositivo político que assegura a soberania epistêmica. Ela garante que a valoração monetária não descaracterize o patrimônio simbólico, convertendo a capacidade absortiva organizacional \cite{CohenLevinthal1990} em poder de barganha real. O processo de externalização do conhecimento (SECI) \cite{Nonaka1995} torna-se, assim, um ato de empoderamento, onde a comunidade retém o controle sobre a definição do que é valioso.

Em última instância, a integração entre machine learning e validação participativa visa superar o hiato entre a riqueza biocultural e a pobreza monetária. A precisão técnica dos indicadores serve à credibilidade institucional; a modelagem econômica viabiliza a inserção em mercados de bioeconomia; e a governança comunitária assegura a justiça distributiva. As adaptações metodológicas, como o uso de proxies e análise de Rasch \cite{Rasch1960}, são estratégias pragmáticas para viabilizar essa tecnologia social em contextos de escassez de dados, estabelecendo um protocolo replicável para que comunidades tradicionais deixem de ser fornecedoras passivas de matéria-prima e assumam a gestão ativa de seu capital intelectual.

Diante desse cenário, esta pesquisa tem como objetivo identificar e analisar o potencial dos sistemas agroecológicos tradicionais como ativos intangíveis, mapeando dimensões de valor cultural, científico e econômico. Partimos da hipótese de que a integração de instrumentos psicométricos e modelos econômicos de valoração, quando submetida à validação comunitária, amplia significativamente a precisão na mensuração do valor estratégico dos conhecimentos tradicionais, transformando saberes difusos em ativos tangíveis para negociação em mercados de valor agregado e fortalecendo a posição negociadora das comunidades quilombolas em redes de inovação colaborativa.

\section{Materiais e métodos}

\subsection{Delineamento e Contexto da Pesquisa}

Esta investigação configura-se como estudo exploratório de métodos mistos \cite{Creswell2018} com abordagem psicométrico-econômica integrada \cite{Messick1995}, desenvolvido junto a comunidades quilombolas do Território de Identidade Semiárido Nordeste II (Bahia), compreendendo dezoito municípios com população total de 425.976 habitantes. A amostra populacional ($n = 384$) foi calculada mediante a fórmula de Cochran \cite{Cochran1977} para populações finitas, considerando margem de erro de 5\%, nível de confiança de 95\% e proporção estimada de 50\%, aplicando-se amostragem estratificada proporcional \cite{Krejcie1970} que considerou a heterogeneidade socioambiental do território. 

Os critérios de inclusão contemplaram produtores agrícolas com idade igual ou superior a 18 anos, de ambos os sexos, residentes nas comunidades certificadas pela Fundação Cultural Palmares e registradas na Base de Informações sobre Povos Indígenas e Localidades Quilombolas do IBGE, com especial ênfase nas onze comunidades quilombolas de Jeremoabo. A coleta de dados foi operacionalizada por equipe de cinco entrevistadores previamente capacitados mediante protocolo padronizado de treinamento \cite{Fowler2014}, incluindo módulos sobre consentimento informado, técnicas de entrevista intercultural e registro sistemático de dados, assegurando uniformidade procedimental, minimização de viés do entrevistador e supervisão sistemática durante todo o processo de campo.

\subsection{Arquitetura do Instrumento Psicométrico}

O instrumento de mensuração foi estruturado em duas dimensões funcionalmente integradas. A primeira dimensão compreendeu variáveis sociodemográficas de caracterização (idade, gênero, escolaridade, ocupação, tempo de residência e envolvimento em atividades agrícolas), permitindo estratificação analítica posterior. A segunda dimensão operacionalizou cinco construtos latentes fundamentados teoricamente: conhecimentos etnoclimatológicos, etnopedológicos e etnoecológicos, representando o domínio cognitivo sobre padrões climáticos, propriedades edáficas e interações ecossistêmicas; percepção do valor estratégico do conhecimento tradicional, mensurando a consciência comunitária sobre a relevância econômica e cultural de seus saberes; capacidade adaptativa às mudanças ambientais, avaliando a plasticidade das práticas frente a variações climáticas e ecológicas; originalidade do conhecimento tradicional, quantificando o grau de singularidade cultural e inovação das práticas; e práticas agroecológicas tradicionais comunitárias, capturando a sofisticação técnica dos sistemas produtivos locais. Cada construto foi mensurado mediante cinco a sete itens formulados em escala Likert de cinco pontos, totalizando trinta itens validados. 

A construção dos itens fundamentou-se na adaptação contextualizada do questionário WOCAT (\textit{World Overview of Conservation Approaches and Technologies}) \cite{Liniger2019} para documentação de tecnologias de Gestão Sustentável da Terra. O instrumento WOCAT, desenvolvido pelo Centre for Development and Environment da Universidade de Berna em parceria com consórcio internacional (ICARDA, SDC, FAO, ISRIC, CIAT, ICIMOD, GIZ), constitui ferramenta padronizada globalmente com validação em mais de 120 países \cite{Schwilch2011} para avaliação sistemática de práticas de manejo sustentável de recursos naturais. 

A arquitetura abrange sete dimensões estruturadas:

\begin{enumerate}[itemsep=0pt,parsep=0pt,topsep=0pt]
\item informações gerais de identificação e localização georreferenciada;
\item descrição técnica detalhada das medidas agronômicas, vegetativas, estruturais e de manejo;
\item classificação tipológica por uso da terra segundo padrões ICRAF;
\item especificações técnicas quantificadas e análise de custos de estabelecimento e manutenção;
\item caracterização multidimensional do ambiente biofísico e socioeconômico mediante indicadores padronizados de clima, topografia, solos, biodiversidade e características dos usuários da terra;
\item avaliação de impactos ambientais, socioeconômicos e produtivos on-site e off-site mediante escala ordinal de sete pontos;
\item análise estruturada de forças, fraquezas, oportunidades de adaptação e vulnerabilidades climáticas percebidas pelos usuários \cite{Liniger2011}.
\end{enumerate} 

A adaptação do questionário WOCAT ao contexto quilombola seguiu protocolo de validação transcultural \cite{Beaton2000}, operacionalizado mediante duas oficinas participativas envolvendo cinco participantes estrategicamente selecionados: dois especialistas em agroecologia, três representantes quilombolas com reconhecida expertise em sistemas tradicionais, e quatro técnicos extensionistas com experiência prévia em pesquisa-ação participativa. 

As oficinas foram estruturadas segundo metodologia de painel de juízes \cite{Alexandre2011}, nas quais os participantes avaliaram independentemente a pertinência cultural, adequação linguística, relevância prática e compreensibilidade de cada item mediante escala de quatro pontos, gerando Índice de Validade de Conteúdo superior a 0,85 para todos os construtos. Este processo iterativo resultou em instrumento híbrido que preserva o rigor metodológico e a comparabilidade internacional do WOCAT enquanto incorpora categorias êmicas \cite{Pike1967} e terminologias vernaculares específicas dos sistemas agroecológicos tradicionais do semiárido brasileiro, assegurando simultaneamente validade de conteúdo transcultural \cite{Milfont2010} e validade ecológica da mensuração \cite{Bronfenbrenner1979}.

\subsection{Validação Psicométrica}

A validação do instrumento seguiu os Standards for Educational and Psychological Testing mediante modelo tripartite. A validade de conteúdo foi aferida por painel de juízes especialistas em agroecologia, psicometria e conhecimentos tradicionais, que avaliaram os itens quanto à clareza, pertinência teórica e relevância prática em escala de 1 a 5 pontos. O Coeficiente de Validade de Conteúdo foi calculado conforme Equação 1, estabelecendo-se $CVC > 0,80$ como critério de aceitação.

\begin{equation}
CVC_i = \frac{k}{j} \times \left(\frac{\sum_{i=1} x_i}{E}\right)
\end{equation}

A validade de estrutura interna integrou Análise Fatorial Exploratória (AFE) \cite{Fabrigar1999} e Confirmatória (AFC) \cite{Brown2015}. Os pressupostos foram verificados mediante teste de esfericidade de Bartlett ($\chi^2$, $p < 0,001$) e índice Kaiser-Meyer-Olkin ($KMO > 0,80$) \cite{Kaiser1974}. A extração fatorial empregou método de Máxima Verossimilhança Robusta (MLR) com rotação oblíqua Promax ($\kappa = 4$) \cite{Satorra1994}, conjugando o critério de Kaiser (autovalores $\lambda > 1,0$), análise paralela de Horn com 1.000 permutações \cite{Horn1965} e inspeção do scree plot \cite{Cattell1966}.

Subsequentemente, a AFC testou o ajuste do modelo pentafatorial mediante bateria de índices: qui-quadrado normalizado $\chi^2/df < 3,0$; Root Mean Square Error of Approximation $RMSEA < 0,06$ (IC90\% $\leq 0,08$) \cite{Browne1992}; Standardized Root Mean Square Residual $SRMR < 0,08$; Comparative Fit Index e Tucker-Lewis Index $CFI$ e $TLI > 0,95$ \cite{Hu1999}; e cargas fatoriais padronizadas $\lambda_{std} > 0,50$ com significância estatística $|z| > 1,96$ ($p < 0,05$) \cite{Hair2019}. 

A invariância psicométrica foi avaliada mediante análise fatorial multigrupo (MGCFA) \cite{Milfont2010}, testando-se sequencialmente invariância configural, invariância métrica ($\Delta CFI < 0,01$) \cite{Cheung2002} e invariância escalar ($\Delta CFI < 0,01$; $\Delta RMSEA < 0,015$) \cite{Chen2007} através do qui-quadrado escalonado Satorra-Bentler ($\Delta\chi^2_{SB}$, $p > 0,05$) \cite{Satorra2001}, assegurando equivalência psicométrica entre subgrupos demográficos.

A confiabilidade foi quantificada mediante tríade de indicadores: alfa de Cronbach ($\alpha > 0,80$) \cite{Cronbach1951}, confiabilidade composta (conforme Equação 2) e variância média extraída (conforme Equação 3) \cite{Fornell1981}. A validade discriminante foi verificada mediante critérios de Fornell-Larcker ($\sqrt{AVE_i} > r_{ij}$) e Heterotrait-Monotrait Ratio ($HTMT < 0,85$) \cite{Henseler2015}.

\begin{equation}
CC = \frac{(\sum\lambda_i)^2}{(\sum\lambda_i)^2 + \sum\theta_i}
\end{equation}

\begin{equation}
AVE = \frac{\sum\lambda_i^2}{\sum\lambda_i^2 + \sum\theta_i}
\end{equation}

\subsection{Modelagem por Teoria de Resposta ao Item}

A validação baseada no padrão de resposta aos itens empregou o Modelo de Crédito Parcial de Masters \cite{Masters1982}, extensão politômica do Modelo Rasch \cite{Rasch1960} apropriada para itens Likert de resposta ordenada. Este modelo permite estimar simultaneamente os limiares de dificuldade ($\delta_{ik}$) entre categorias adjacentes e a habilidade latente ($\theta_n$) dos respondentes na escala logit, conforme Equação 12 \cite{Bond2020}.

\begin{equation}
P(X_{ni}=x|\theta_n, \delta_i) = \frac{\exp\left(\sum_{k=0}^{x}(\theta_n - \delta_{ik})\right)}{\sum_{j=0}^{m_i}\exp\left(\sum_{k=0}^{j}(\theta_n - \delta_{ik})\right)}
\end{equation} 

O ajuste item-pessoa foi avaliado mediante índices de adequação Infit e Outfit MNSQ, acompanhados de suas respectivas estatísticas padronizadas (ZSTD) \cite{Wright1994}. Estabeleceram-se como critérios de aceitação: valores de MNSQ entre 0,6-1,4 logits \cite{Linacre2002}; ZSTD $< \pm 2,0$ ($p < 0,05$); e correlação ponto-medida superior a 0,30 \cite{Bond2020}. Itens que apresentaram desajuste severo (MNSQ > 1,5 ou ZSTD > 3,0) foram inspecionados qualitativamente mediante análise de conteúdo estruturada, consulta ao painel de juízes e triangulação com curvas características do item (ICC), sendo então reformulados ou eliminados quando necessário \cite{Wilson2005}.

A confiabilidade da separação foi quantificada através dos coeficientes de separação pessoa (conforme Equação 4) e item ($G_{item}$) \cite{Wright1996}, estabelecendo-se limiares de $G_{person} > 0,80$ e $G_{item} > 0,90$ \cite{Fisher2007}, complementados pelos índices Person Reliability Index (PRI) e Item Reliability Index (IRI) \cite{Bond2020}.

\begin{equation}
G_{person} = \frac{SD_{adjusted}^2}{SD_{adjusted}^2 + SE_{mean}^2}
\end{equation}

O funcionamento diferencial dos itens (DIF) foi avaliado mediante procedimento de Mantel-Haenszel \cite{Mantel1959} com correção de Bonferroni para comparações múltiplas. A magnitude do DIF foi quantificada através do contraste em logits (conforme Equação 5), interpretado segundo a classificação ETS: $|DIF| < 0,43$ logits (negligível); $0,43 \leq |DIF| < 0,64$ (leve a moderado); $|DIF| \geq 0,64$ (severo) \cite{Zieky1993}. Adicionalmente, o DIF não-uniforme foi testado mediante análise de regressão logística ordinal \cite{Swaminathan1990}, considerando-se $\Delta R^2 > 0,035$ como indicativo de efeito substancial \cite{Jodoin2003}.

\begin{equation}
DIF_{contrast} = \delta_{focal} - \delta_{reference}
\end{equation}

As análises foram conduzidas mediante o software Winsteps versão 5.5.2 \cite{Linacre2023}, que gerou mapas pessoa-item (Wright maps) \cite{Wilson2005} para inspeção visual simultânea da adequação da cobertura do traço latente, identificação de intervalos de medida (measurement gaps) superiores a 0,5 logit, detecção de efeitos de teto ou piso, e avaliação do targeting apropriado do instrumento ($|\theta_{mean} - \delta_{mean}| < 1$ logit) \cite{Bond2020}.

\subsection{Integração com Modelos de Valoração Econômica}

Os escores psicométricos obtidos via Teoria de Resposta ao Item foram integrados ao framework Value Added Intellectual Coefficient (VAIC) \cite{Pulic2000} para quantificação do valor agregado pelos conhecimentos tradicionais. O modelo VAIC, originalmente desenvolvido para mensuração de capital intelectual corporativo, foi adaptado ao contexto comunitário mediante decomposição tricotômica de componentes \cite{Bontis2000}: (1) eficiência do capital humano (HCE), calculada conforme Equação 6, refletindo a produtividade do conhecimento incorporado nos indivíduos \cite{Edvinsson1997}; (2) eficiência do capital estrutural (SCE), mensurada conforme Equação 7, capturando a infraestrutura organizacional de conhecimento \cite{Stewart1997}; e (3) eficiência do capital empregado (CEE), expressa conforme Equação 8, representando a produtividade dos ativos físicos \cite{Roos1997}.

\begin{equation}
HCE = \frac{VA}{HC}
\end{equation}

\begin{equation}
SCE = \frac{SC}{VA} = \frac{VA - HC}{VA}
\end{equation}

\begin{equation}
CEE = \frac{VA}{CE}
\end{equation}

Os escores de proficiência em conhecimento tradicional ($\theta_n$) derivados da modelagem Rasch foram empregados como proxies para quantificação do capital intelectual comunitário \cite{Sveiby1997}, estabelecendo-se a correspondência formal entre os construtos psicométricos e os componentes do capital humano conforme Equação 9, onde os coeficientes $\alpha$ e $\beta$ foram calibrados por regressão múltipla com dados de produtividade agrícola como variável dependente \cite{Becker1993}. O cálculo do VAIC total seguiu a formulação aditiva conforme Equação 10 \cite{Pulic2000}, onde valores superiores a 3,0 indicam criação eficiente de valor \cite{Iazzolino2013}, e a decomposição percentual de cada componente permitiu identificar a contribuição relativa dos ativos tangíveis versus intangíveis para a sustentabilidade dos sistemas agroecológicos tradicionais.

\begin{equation}
HC = \alpha + \beta\theta_n
\end{equation}

\begin{equation}
VAIC = CEE + HCE + SCE
\end{equation}

A valoração econômica incorporou adicionalmente a metodologia de custo de reposição \cite{Freeman2003}, estimando o investimento hipotético necessário para reconstituição de conhecimentos perdidos conforme Equação 11, onde $C_i$ representa o custo unitário de capacitação, $T_i$ o tempo de aprendizagem estimado e $W_i$ os ponderadores de complexidade e raridade derivados da análise psicométrica \cite{Bateman2002}. A integração metodológica triangulada permitiu a tradução de atributos qualitativos em métricas quantitativas auditáveis \cite{Kaplan2004}, fornecendo base probatória rigorosa para negociações de transferência de tecnologia, contratos de repartição de benefícios sob o Protocolo de Nagoia \cite{Laird2020} e certificação de produtos com indicação geográfica protegida \cite{Belletti2015}.

\begin{equation}
CRP = \sum_{i=1}^{n} C_i \times T_i \times W_i
\end{equation}

\subsection{Modelagem Fuzzy e Construção do Índice de Valoração Bioeconômica (IVB)}

O Índice de Valoração Bioeconômica (IVB) será construído mediante sistema de inferência fuzzy tipo Mamdani \cite{Mamdani1975}, integrando as variáveis linguísticas validadas pelo painel Delphi (Capítulo~2) com os escores psicométricos TRI e os indicadores econômicos VAIC. A arquitetura do sistema compreende três módulos integrados:

\textbf{Módulo 1 --- Fuzzificação.} As variáveis de entrada serão fuzzificadas mediante funções de pertinência triangulares e trapezoidais calibradas com base nos limiares de consenso do Delphi (mediana, IQR) e nos parâmetros TRI ($\theta_n$, $\delta_{ik}$). Para cada variável linguística validada, três a cinco termos linguísticos (e.g., \textit{muito baixo}, \textit{baixo}, \textit{médio}, \textit{alto}, \textit{muito alto}) serão definidos com universos de discurso ancorados nas distribuições empíricas dos respondentes e nos consensos dos especialistas. Os escores VAIC (HCE, SCE, CEE) serão fuzzificados como variáveis de entrada adicionais, representando a dimensão econômica da valoração.

\textbf{Módulo 2 --- Base de regras SE-ENTÃO.} As regras de inferência serão construídas mediante elicitação híbrida: (i)~regras derivadas diretamente do consenso Delphi, onde as relações de importância entre variáveis validadas ($CVC \geq 0,80$; $W \geq 0,70$) informam as proposições condicionais; (ii)~regras derivadas dos modelos TRI, onde os limiares de dificuldade e os padrões de resposta ao item indicam interações não-lineares entre dimensões; e (iii)~regras derivadas da validação participativa, incorporando relações causais identificadas pelos mestres de saberes nas oficinas comunitárias. A consistência da base de regras será verificada computacionalmente para eliminação de contradições e redundâncias \cite{Zadeh1965}.

\textbf{Módulo 3 --- Defuzzificação e IVB.} O método do centroide será empregado para defuzzificação, produzindo um índice numérico contínuo (IVB $\in [0, 100]$) que sintetiza as múltiplas dimensões de valoração. O IVB será decomposto em sub-índices correspondentes às dimensões validadas (cultural-simbólica, biofísica-ambiental, econômica-mercadológica, institucional-governança, adaptativa-resiliência, social-organizacional), permitindo análise diferenciada por componente e comparação inter-comunitária.

A validação do sistema fuzzy seguirá protocolo dual: (i)~validação técnica mediante análise de sensibilidade, superfícies de resposta e comparação com rankings obtidos pela agregação linear ponderada dos escores VAIC; e (ii)~validação participativa mediante oficinas devolutivas onde representantes comunitários avaliarão a congruência entre o IVB computado e suas percepções sobre o valor relativo dos sistemas avaliados. Discrepâncias significativas entre valoração computacional e percepção comunitária acionarão ciclos de revisão da base de regras, assegurando conformidade com o princípio de validade consequencial \cite{Messick1995}.


\subsection{Validação Participativa e Validade Consequencial}

A validade consequencial foi estabelecida mediante delineamento adaptado do framework interpretativo de Kane \cite{Kane2013}, operacionalizando quatro componentes integrados do argumento de validação \cite{Kane2006}: (1) especificação dos efeitos pretendidos do sistema de avaliação sobre empoderamento econômico e fortalecimento da soberania epistêmica comunitária \cite{Messick1995}; (2) descrição explícita dos componentes do instrumento com fundamentação teórica na Teoria dos Recursos da Firma \cite{Barney1991} e na Economia Ecológica \cite{Costanza1997}; (3) formulação de afirmações interpretativas derivadas dos resultados, conectando escores psicométricos a capacidades de gestão de conhecimento \cite{Shavelson2009}; e (4) delineamento de mecanismos de ação designados para produção dos efeitos esperados mediante fortalecimento da capacidade absortiva organizacional \cite{CohenLevinthal1990}.

A validação participativa foi operacionalizada através de quatro oficinas comunitárias estruturadas segundo o método de grupos focais \cite{Krueger2014}, nas quais os resultados preliminares da análise psicométrica foram apresentados mediante visualizações de dados acessíveis (Wright maps, gráficos radar de competências) e submetidos à interpretação crítica de 28 participantes quilombolas, incluindo agricultores experientes, lideranças comunitárias e jovens em processo de aprendizagem \cite{FalsBorda1991}. Cada oficina seguiu protocolo estruturado de validação semântica \cite{Pasquali2010}, verificando-se a congruência entre as métricas quantitativas e as percepções qualitativas dos detentores dos conhecimentos mediante análise temática de conteúdo \cite{Braun2006}.

Este processo iterativo de validação comunicativa \cite{Habermas1984} garantiu que a valoração econômica não se descolasse das dimensões simbólicas e relacionais que conferem significado aos saberes tradicionais \cite{Agrawal2002}, prevenindo mercantilização reducionista \cite{Polanyi1944} e apropriação indébita mediante salvaguardas jurídicas fundamentadas no consentimento livre, prévio e informado \cite{Wynberg2015}. A triangulação metodológica \cite{Denzin1978} entre dados quantitativos psicométricos, valoração econômica VAIC e validação qualitativa participativa produziu síntese validada que fundamenta tanto a teoria da mensuração de ativos intangíveis \cite{Lev2001} quanto as estratégias práticas de gestão de conhecimentos tradicionais em contextos de justiça cognitiva \cite{Santos2007} e repartição equitativa de benefícios \cite{Schroeder2020}.







%Textos para discussão



%A operacionalização dessa integração metodológica enfrenta, entretanto, um desafio epistemológico crucial relacionado à validação comunitária dos instrumentos de mensuração. A aplicação acrítica de escalas psicométricas desenvolvidas sob paradigmas positivistas ocidentais carrega o risco de violência epistêmica descrita por Santos \cite{Santos2007}, impondo categorias analíticas exógenas que distorcem ou invisibilizam as ontologias nativas de conhecimento. A teoria da validade consequencial proposta por Messick \cite{Messick1995} alerta que a validade de um instrumento de mensuração não se restringe à sua precisão técnica, mas inclui as consequências sociais de sua aplicação, exigindo que os processos de valoração sejam submetidos à validação comunitária mediante protocolos participativos que assegurem a congruência cultural dos indicadores. Essa exigência metodológica alinha-se ao conceito de pesquisa-ação participativa de Fals Borda \cite{FalsBorda1991}, onde os sujeitos da pesquisa atuam como co-pesquisadores na definição dos parâmetros de valoração, garantindo que as dimensões de valor cultural, simbólico e territorial reconhecidas pelas comunidades sejam incorporadas aos modelos econômicos, transcendendo as limitações reducionistas de abordagens puramente monetárias.
%A validação comunitária dos instrumentos de valoração cumpre uma função estratégica adicional ao fortalecer a capacidade absortiva organizacional das comunidades em relação aos seus próprios ativos intangíveis. Cohen e Levinthal \cite{CohenLevinthal1990} definem capacidade absortiva como a habilidade de uma organização em reconhecer o valor de informações externas, assimilá-las e aplicá-las para fins comerciais, conceito que pode ser transposto para o contexto comunitário como a capacidade de identificar, valorar e mobilizar estrategicamente os conhecimentos tradicionais como ativos negociáveis. O processo participativo de construção e validação de métricas psicométrico-econômicas atua como mecanismo de externalização do conhecimento tácito na espiral SECI proposta por Nonaka e Takeuchi \cite{Nonaka1995}, convertendo saberes implícitos em conhecimento explícito codificado sem comprometer sua autenticidade cultural. Essa externalização instrumentalizada constitui prerequisito para que as comunidades assumam protagonismo nas negociações de repartição de benefícios derivados da biodiversidade, licenciamento de tecnologias sociais e certificação de produtos tradicionais, superando a posição histórica de passividade informacional que caracteriza as relações assimétricas com agentes externos.

%A integração entre psicometria, modelagem econômica e validação comunitária configura, portanto, um framework metodológico triádico capaz de superar a invisibilidade econômica dos conhecimentos tradicionais. A precisão psicométrica garante a robustez científica da mensuração, conferindo credibilidade aos indicadores de proficiência comunitária perante agências de financiamento, instituições de pesquisa e mercados de valor agregado. A modelagem econômica através do VAIC traduz essas métricas em linguagem monetária inteligível aos sistemas de contabilidade corporativa e avaliação de impacto, permitindo a incorporação dos ativos intangíveis tradicionais nos cálculos de valoração territorial e nas estratégias de desenvolvimento local. A validação comunitária, por sua vez, assegura a legitimidade cultural dos processos de mensuração e garante que a valoração econômica não se descole das dimensões simbólicas e relacionais que conferem significado aos saberes tradicionais, prevenindo a mercantilização reducionista e a apropriação indébita por atores externos.

%Contudo, a efetivação desse framework exige a superação de obstáculos metodológicos concretos relacionados à operacionalização dos instrumentos psicométricos em contextos de alta diversidade cultural e baixa escolaridade formal. A aplicação de escalas TRI demanda amostras robustas e múltiplos itens calibrados para diferentes níveis de dificuldade, o que implica em processos longos de coleta de dados e análise estatística que frequentemente excedem os recursos disponíveis em projetos de pesquisa-ação participativa. Adicionalmente, a modelagem do VAIC pressupõe a disponibilidade de dados contábeis estruturados sobre outputs produtivos e inputs de capital, informações raramente sistematizadas em economias comunitárias baseadas em trocas não-monetárias e arranjos de reciprocidade. Essas limitações práticas justificam a necessidade de adaptações metodológicas que mantenham o rigor psicométrico e econômico enquanto simplificam os procedimentos operacionais, mediante estratégias como amostragem adaptativa, escalas reduzidas validadas por análise de Rasch \cite{Rasch1960} e proxies econômicas baseadas em preços hedônicos e valoração contingente adaptada culturalmente.

%A integração entre psicometria, modelagem econômica e validação comunitária configura, portanto, um framework metodológico triádico capaz de superar a invisibilidade econômica dos conhecimentos tradicionais. A precisão psicométrica garante a robustez científica da mensuração, conferindo credibilidade aos indicadores de proficiência comunitária perante agências de financiamento, instituições de pesquisa e mercados de valor agregado. A modelagem econômica através do VAIC traduz essas métricas em linguagem monetária inteligível aos sistemas de contabilidade corporativa e avaliação de impacto, permitindo a incorporação dos ativos intangíveis tradicionais nos cálculos de valoração territorial e nas estratégias de desenvolvimento local. A validação comunitária, por sua vez, assegura a legitimidade cultural dos processos de mensuração e garante que a valoração econômica não se descole das dimensões simbólicas e relacionais que conferem significado aos saberes tradicionais, prevenindo a mercantilização reducionista e a apropriação indébita por atores externos.

%A contribuição teórica desta pesquisa reside na construção de uma ponte metodológica entre três tradições disciplinares historicamente desconexas: a psicometria educacional, a economia de ativos intangíveis e a pesquisa participativa em contextos de comunidades tradicionais. Ao demonstrar a viabilidade técnica e a validade cultural dessa integração, o estudo fornece um protocolo replicável para processos de valoração de conhecimentos tradicionais que respeitem simultaneamente os critérios de rigor científico demandados pela academia e pelos mercados de conhecimento, bem como os princípios de autodeterminação informacional e soberania epistemológica reivindicados pelos movimentos sociais de povos e comunidades tradicionais. A contribuição prática manifesta-se na instrumentalização das comunidades quilombolas com métricas auditáveis de valoração de seus ativos intangíveis, capacitando-as para negociações mais simétricas em processos de repartição de benefícios, licenciamento de tecnologias sociais e acesso a mercados diferenciados que remuneram atributos de qualidade e territorialidade.

%A urgência dessa investigação justifica-se pela conjuntura crítica de erosão acelerada dos sistemas tradicionais sob pressão das mudanças climáticas e da expansão de monoculturas industriais, simultânea à crescente demanda por soluções baseadas na natureza e conhecimentos de adaptação climática que as comunidades tradicionais detêm de forma privilegiada. A ausência de métricas adequadas de valoração mantém esses conhecimentos na condição de externalidades positivas não remuneradas, perpetuando a vulnerabilidade econômica das comunidades e incentivando o abandono geracional de práticas agroecológicas sofisticadas. A construção de frameworks robustos de mensuração e valoração converte essa dinâmica perversa ao posicionar os saberes tradicionais como ativos estratégicos remuneráveis, criando incentivos econômicos endógenos para sua conservação e transmissão intergeracional, alinhados aos princípios de justiça cognitiva e repartição equitativa de benefícios preconizados pelo Protocolo de Nagoya e pela Convenção da Diversidade Biológica.

%Diante deste cenário, esta pesquisa tem como objetivo identificar e analisar o potencial dos sistemas agroecológicos tradicionais como ativos intangíveis, mapeando dimensões de valor cultural, científico e econômico que orientem sua valoração comunitária e incorporação em estratégias de gestão do conhecimento. Partimos da hipótese de que a integração de instrumentos psicométricos e modelos econômicos de valoração, quando submetida à validação comunitária, amplia significativamente a precisão na mensuração do valor estratégico dos conhecimentos tradicionais, transformando saberes difusos em ativos tangíveis para negociação em mercados de valor agregado e fortalecendo a posição negociadora das comunidades quilombolas em redes de inovação colaborativa. 


















